\documentclass[12pt,a4paper]{article}
\usepackage[margin=25mm, headheight=15pt, headsep=10pt]{geometry}
\usepackage{fontspec}
\usepackage[table]{xcolor} % Note: table option is required for rowcolors
\usepackage{titlesec}
\usepackage{tocloft}
\usepackage{fancyhdr}
\usepackage{booktabs}
\usepackage{tcolorbox}
\tcbuselibrary{skins, breakable}
\usepackage[colorlinks=true, linkcolor=emerald, urlcolor=emerald, bookmarks=true]{hyperref}

\definecolor{emerald}{RGB}{0,102,51}
\definecolor{darkred}{RGB}{139,0,0}
\definecolor{lightgray}{RGB}{245,245,245}

\usepackage[nil]{babel}
\babelprovide[import, main]{bengali}
\babelprovide[import]{arabic}
\babelfont{rm}[Renderer=HarfBuzz,Scale=1.0]{Noto Serif Bengali}
\babelfont{sf}[Renderer=HarfBuzz]{Noto Sans Bengali}
\babelfont{tt}[Renderer=HarfBuzz]{Noto Sans Bengali}
\babelfont[arabic]{rm}[Renderer=HarfBuzz,Scale=1.1]{Noto Naskh Arabic}

% Section styling
\titleformat{\section}{\Large\bfseries\color{emerald}}{\thesection.}{0.5em}{}
\titleformat{\subsection}{\large\bfseries\color{emerald!85!black}}{\thesubsection.}{0.5em}{}
\titleformat{\subsubsection}{\normalsize\bfseries\color{emerald!70!black}}{\thesubsubsection.}{0.5em}{}
\titlespacing*{\section}{0pt}{18pt}{8pt}

% Fancy Header/Footer setup
\pagestyle{fancy}
\fancyhf{} % clear all header and footer fields
\fancyhead[L]{\color{emerald}\small\textbf{\nouppercase{\leftmark}}}
\fancyfoot[L]{\scriptsize\color{black!50}{\lat{AI}}-সংকলিত, আলেম-যাচাইকৃত নয়}
\fancyfoot[C]{\color{emerald!70!black}\thepage}
\renewcommand{\headrulewidth}{0.4pt}
\renewcommand{\headrule}{\hbox to\headwidth{\color{emerald}\leaders\hrule height \headrulewidth\hfill}}
\renewcommand{\footrulewidth}{0pt}

% Custom boxes
% 1. Ayat/Hadith Box
\newtcolorbox{ayatbox}[1][]{
    enhanced, breakable, 
    colback=emerald!5, 
    colframe=emerald, 
    boxrule=0.5pt, 
    arc=3pt, 
    left=10pt, right=10pt, top=10pt, bottom=10pt,
    #1
}

% 2. Quote/Translation Box
\newtcolorbox{quotebox}[1][]{
    enhanced, breakable,
    frame hidden,
    colback=lightgray,
    borderline west={3pt}{0pt}{emerald},
    left=15pt, right=10pt, top=10pt, bottom=10pt,
    sharp corners,
    #1
}

% 3. AI-generated notice box
\newtcolorbox{warnbox}[1][]{
    enhanced, breakable,
    colback=red!4,
    colframe=darkred,
    boxrule=0.8pt,
    arc=3pt,
    left=10pt, right=10pt, top=10pt, bottom=10pt,
    #1
}

% Commands
\newcommand{\arab}[1]{\foreignlanguage{arabic}{#1}}
\newcommand{\kib}[1]{\textcolor{emerald}{\textbf{#1}}}

% Latin (English) fallback: Noto Serif Bengali has NO Latin glyphs
\newfontfamily\latfont[Renderer=HarfBuzz]{Noto Serif}
\newcommand{\lat}[1]{{\latfont #1}}

\setlength{\parskip}{5pt}
\setlength{\parindent}{0pt}

% Fix for missing bullet in Noto Serif Bengali
\renewcommand{\labelitemi}{$\bullet$}



\begin{document}
\begin{center}{\Large\bfseries\color{emerald} ঘটনা ১ — আয়েশা (রাঃ): নবীজির কুশন ও ছবিওয়ালা পর্দা}\end{center}
\vspace{1em}
\section*{ঘটনা ১ — আয়েশা (রাঃ): নবীজির কুশন ও ছবিওয়ালা পর্দা}

\begin{warnbox}
 \textbf{গুরুত্বপূর্ণ নোটিশ (\lat{Important} \lat{Notice}):} এই কন্টেন্টটি \lat{AI} (কৃত্রিম বুদ্ধিমত্তা) দিয়ে প্রস্তুত ও সংকলিত। \lat{AI} কঠোর সূত্র-মান নীতি মেনে শুধু নির্ভরযোগ্য উৎস থেকে তথ্য সংগ্রহ করেছে — কেবল সহীহ/হাসান হাদিস ও সহীহ সনদের আসার রাখা হয়েছে, দুর্বল সনদ বাদ দেওয়া হয়েছে। তবে এটি কোনো আলেম বা মুহাদ্দিস কর্তৃক যাচাইকৃত নয়।
\end{warnbox}

\begin{quote}
\textbf{সম্পর্কিত ফাইল:} এই দ্বিধার হাদিস ও ব্যবহারিক গাইড — ছবির বিধানের সবচেয়ে গুরুত্বপূর্ণ দুইটি সাহাবি ঘটনা এখানে।
\end{quote}

\medskip\hrule\medskip

\subsection*{১. সূত্র কার্ড (\lat{Source} \lat{Card})}

\begin{table}[h]\centering\small
\begin{tabular}{ll}
বিষয় & বিবরণ \\
\hline
\textbf{ঘটনা} & আয়েশা (রাঃ) ছবিওয়ালা কুশন কিনলেন $\rightarrow$ নবীজি (সা.) ঘরে ঢুকলেন না $\rightarrow$ পরে পর্দার ছবি টেনে ছিঁড়লেন $\rightarrow$ ছবিকে বালিশ বানালে আপত্তি নেই \\
\textbf{রাবি} & আয়েশা (রাঃ) — নবীজির স্ত্রী \\
\textbf{গ্রন্থ ও নম্বর} & সহীহ বুখারি ৫৯৫৭, ৫৯৬১, ৫৯৫৪; সহীহ মুসলিম ২১০৭, ২১০৮ \\
\textbf{সনদের মান} & \textbf{সহীহ} (বুখারি-মুসলিম) — সর্বোচ্চ মানের সনদ \\
\textbf{স্থান ও সময়} & মদিনা — হিজরত-পরবর্তী যুগ, নবীজির (সা.) জীবদ্দশায় \\
\hline
\end{tabular}
\end{table}

\medskip\hrule\medskip

\subsection*{২. পটভূমি (\lat{Background})}

\textbf{কে:} আয়েশা বিনতে আবু বকর (রাঃ) — নবীজির (সা.) প্রিয় স্ত্রী ও তালিমপ্রাপ্ত মুমিনা; উম্মুল মুমিনীন।

\textbf{কখন ও কোথায়:} মদিনায়, নবীজির (সা.) সংসারে — দুটি ভিন্ন সময়ের দুইটি পর্ব: (১) কুশন কেনার ঘটনা, (২) পর্দা লাগানোর ঘটনা।

\textbf{কেন:} ঘর সাজানোর স্বাভাবিক \textbf{ইচ্ছা (\lat{desire})} থেকে আয়েশা (রাঃ) ছবিওয়ালা কুশন কিনলেন, এবং পরে দরজায় ছবিওয়ালা পর্দা লাগালেন — নবীজি (সা.) দুইবারই প্রতিক্রিয়া দেখিয়ে ছবি নিয়ে ইসলামের বিধান চিরতরে স্থির করে দিলেন। এই দুইটি ঘটনাই ছবির বিধানের সবচেয়ে গুরুত্বপূর্ণ ভিত্তি — নিষেধের কঠোরতা এবং ব্যতিক্রমের সীমারেখা দুই-ই এখানে।

\medskip\hrule\medskip

\subsection*{৩. পূর্ণ ঘটনার বর্ণনা (\lat{The} \lat{Full} \lat{Event})}

\subsubsection*{পর্ব ১ — কুশন কেনার ঘটনা}

আয়েশা (রাঃ) বাজারে একটি \textbf{কুশন (নুমরুকা)} দেখলেন — যার গায়ে ছবি ছিল। তিনি সেটি কিনলেন; নিয়ত ছিল — নবীজি (সা.) এতে বসবেন ও হেলান দেবেন। ঘরে এনে রাখলেন।

নবীজি (সা.) কুশনটি দেখে দরজার কাছে এসে \textbf{থেমে গেলেন — ভেতরে প্রবেশ করলেন না।} আয়েশা (রাঃ) নবীজির চেহারায় অপ্রসন্নতার লক্ষণ দেখে চমকে গেলেন। তিনি সঙ্গে সঙ্গে বললেন:

\begin{quote}
\lat{“}হে আল্লাহর রাসূল! আমি আল্লাহ ও তাঁর রাসূলের কাছে তাওবা করছি — আমি কী অপরাধ করেছি?\lat{”}
\end{quote}

নবীজি (সা.) প্রশ্ন করলেন:

\begin{quote}
\lat{“}এই কুশনের ব্যাপারটি কী?\lat{”}
\end{quote}

আয়েশা (রাঃ) বললেন — আপনার বসা ও হেলান দেওয়ার জন্য কিনেছি। নবীজি (সা.) তখন বললেন:

\begin{quote}
\lat{“}এই ছবিগুলোর মালিকদের কিয়ামতের দিন শাস্তি দেওয়া হবে, আর বলা হবে — যা সৃষ্টি করেছ, তাতে প্রাণ দাও।\lat{”} আরও বললেন: \textbf{\lat{“}নিশ্চয়ই যে ঘরে ছবি থাকে, সে ঘরে ফেরেশতারা প্রবেশ করেন না।\lat{”}}
\end{quote}

\subsubsection*{পর্ব ২ — পর্দার ঘটনা}

আরেকবার আয়েশা (রাঃ) দরজায় একটি \textbf{পর্দা (কিরাম)} লাগালেন — যার গায়েও ছবি ছিল। নবীজি (সা.) এক সফর থেকে ফিরে ঘরে এসে পর্দাটি দেখলেন — \textbf{চেহারার রং বদলে গেল (রাগের আলামতে)।} তিনি পর্দাটি টেনে নামালেন, এমনকি টেনে ছিঁড়ে ফেললেন অথবা টুকরা করে দিলেন, আর বললেন:

\begin{quote}
\lat{“}আল্লাহ আমাদেরকে পাথর বা মাটিকে পোশাক পরানোর নির্দেশ দেননি।\lat{”}
\end{quote}

অন্য বর্ণনায় তিনি বললেন:

\begin{quote}
\lat{“}আয়েশা! কিয়ামতের দিন মানুষের মধ্যে সবচেয়ে কঠিন শাস্তি হবে তাদের, যারা আল্লাহর সৃষ্টির \textbf{অনুকরণ (\lat{imitation})} করে।\lat{”}
\end{quote}

আয়েশা (রাঃ) বিষণ্ণ হলেন, কিন্তু নবীজির (সা.) শিক্ষা মেনে তিনি পর্দাটির \textbf{কাপড় কেটে দুইটি বালিশ} বানালেন এবং ভেতরে খেজুরের আঁশ ভরে দিলেন। এরপর নবীজি (সা.) সেই বালিশে \textbf{হেলান দিতে} শুরু করলেন — কোনো আপত্তি করলেন না।

\medskip\hrule\medskip

\subsection*{৪. শব্দের ব্যাখ্যা (\lat{Key} \lat{Words})}

\begin{itemize}
\item \textbf{নুমরুকাহ (\arab{نمرقة})}: কুশন/গদি — বসা ও হেলান দেওয়ার বালিশ।
\item \textbf{তাসাবির (\arab{تصاوير})}: ছবি — প্রাণবান সৃষ্টির ছবি।
\item \textbf{কিরাম (\arab{قرام})}: পর্দা — দরজায় ঝোলানো ছবিওয়ালা চাদর।
\item \textbf{মুমতাহানাহ (\arab{ممتهنة})}: অবজ্ঞার/ব্যবহৃত বস্তু — মাটিতে মোছা, বসা, হেলান দেওয়া হয় এমন জিনিস।
\item \textbf{আল-মুসাওয়িরুন (\arab{المصورون})}: ছবি নির্মাতারা — প্রাণবান সৃষ্টির নকলকারী।
\item \textbf{উসাদাতান (\arab{وسادتان})}: দুইটি বালিশ — পর্দার কাপড় থেকে তৈরি।
\end{itemize}
\medskip\hrule\medskip

\subsection*{৫. সংশ্লিষ্ট হাদিস ও আয়াত}

\begin{quote}
\textbf{\lat{“}যে ঘরে ছবি বা কুকুর থাকে, ফেরেশতারা সে ঘরে প্রবেশ করেন না।\lat{”}} — সহীহ বুখারি ৩২২৬, সহীহ মুসলিম ২১০৬
\end{quote}

\begin{quote}
\textbf{\lat{“}যদি ছবি আঁকতেই চাও, তাহলে গাছের ছবি আঁকো — এমন বস্তুর ছবি আঁকো যাতে প্রাণ নেই।\lat{”}} — ইবনে আব্বাস (রাঃ), সহীহ মুসলিম ২১১০খ
\end{quote}

\begin{quote}
\textbf{\lat{“}যারা আল্লাহকে বাদ দিয়ে যাদের ডাকে, তারা একটি মাছিও সৃষ্টি করতে পারবে না।\lat{”}} — সূরা আল-হজ্জ ২২:৭৩ (বাংলা অর্থ)
\end{quote}

\begin{quote}
\textbf{\lat{“}নিশ্চয়ই আল্লাহ পাক-পবিত্র মানুষকে পছন্দ করেন...\lat{”}} — প্রসঙ্গ: ছবির নিষেধ মূলত অন্তরের পবিত্রতা ও তাওহীদের সংরক্ষণ (সূরা আত-তাওবাহ ৯:১০৮-এর ভাবনার আলোকে)।
\end{quote}

\medskip\hrule\medskip

\subsection*{৬. রেফারেন্স}

\begin{itemize}
\item সহীহ বুখারি: ৫৯৫৭ (কুশন), ৫৯৬১ (কুশন — পূর্ণ বর্ণনা), ৫৯৫৪ (পর্দা — বালিশ রূপান্তর), ৩২২৬ (ফেরেশতা-কুকুর-ছবি)
\item সহীহ মুসলিম: ২১০৭ (কুশন), ২১০৮ (পর্দা-বালিশ)
\item সহীহ মুসলিম: ২১১০খ (ইবনে আব্বাস — গাছের ছবি)
\item ব্যাখ্যা: শারহু সহীহ মুসলিম (ইমাম নববি), ফাতহুল বারি (ইবনে হাজার)
\end{itemize}

\end{document}
